\subsection{重心的定义}

设 E 是 R 上的一个 Hausdorff 局部凸空间，\(E'\) 是它的对偶，\(E'^{*}\) 是 \(E'\) 的代数对偶，E 被规范地等同于 \(E'^{*}\) 的一个线性子空间。设 K 是 E 的一个紧子集；K 到 E 中的规范单射连续且具有紧支集，故对 K 上的每个测度 \(\mu\)，积分 \(\int x d\mu(x)\) 有定义且是 \(E'^{*}\) 的一个元素（Ch. III, §3, No.~1）。此外，在 K 上，由弱拓扑 \(\sigma(E'^{*}, E')\) 诱导的拓扑与原来的拓扑一致。最后，若 C 是 K 在 \(E'^{*}\)（赋以 \(\sigma(E'^{*}, E')\)）中的闭凸包，则 \(C \cap E\) 是 K 在 E 中关于原来的拓扑（或关于弱化拓扑 \(\sigma(E, E')\)）的闭凸包。

定义 1. --- 设 K 是 Hausdorff 局部凸空间 E 的一个紧子集。对 K 上每个总质量等于 1 的正测度 \(\mu\)，向量 \(\mathbf{b}_{\mu} = \int \mathbf{x} d\mu(\mathbf{x})\)（属于 \(E'^{*}\)）称为 \(\mu\) 的重心。

例. --- 设 \(\mu\) 是 K 上的一个离散测度，为正且总质量为 1；于是它形如 \(\mu = \sum_{i=1}^{n} \lambda_i \varepsilon_{x_i}\)，其中 \(x_i \in K\)，诸 \(\lambda_i\) 是实数，满足对一切 i 有 \(\lambda_i \geqslant 0\) 且 \(\sum_{i=1}^{n} \lambda_i = 1\) 。由于 \(\int x d\varepsilon_y(x) = y\)（Ch. III, §3, No.~1, 例 3），\(b_\mu = \int x d\mu(x) = \sum_{i} \lambda_i x_i\) 。特别地，对每个 \(x \in K\)，x 是测度 \(\varepsilon_x\) 的重心。

命题 1. --- 设 E 是一个 Hausdorff 局部凸空间，K 是 E 的一个紧子集，C 是 K 在 E 中的闭凸包。则 C 由 E 中至少是 K 上一个质量为 1 的正测度的重心的那些点组成。

这不过是 Ch. III, §3, No.~2 命题 5 应用于 K 到 E 中的规范单射而已。

推论. --- 若 K 在 E 中的闭凸包 C 是紧的，则 K 上每个总质量为 1 的正测度的重心属于 E。

因为，这时 C 也是 K 在赋以弱拓扑 \(\sigma (\mathrm{E}^{\prime *},\mathrm{E}^{\prime})\) 的 E'* 中的闭凸包，故只需对 K 到 E 中的规范单射应用 Ch. III, §3, No.~2 命题 4。

注. --- 命题 1 的推论特别适用于 K 为凸集或 E 为拟完备的情形。

命题 2. --- 设 K 是 Hausdorff 局部凸空间 E 的一个紧凸子集，\(\mu\) 是 K 上总质量为 1 的一个正测度，\(b_{\mu}\) 是它的重心。对每个在 K 上下半连续的凸数值函数 \(f \geqslant 0\)，

\[
f (\mathbf {b} _ {\mu}) \leqslant \int^ {*} f d \mu .
\]

已知（TVS, II, §5, No.~4, 命题 5）\(f\) 是到 \(K\) 上的限制——连续仿射线性函数 \(h_{\alpha} : \mathbf{x} \mapsto c_{\alpha} + \langle \mathbf{x}, \mathbf{z}_{\alpha}' \rangle\) 的限制——所成族的上包络。则

\[
\int h _ {\alpha} (\mathbf {x}) d \mu (\mathbf {x}) \leqslant \int^ {*} f (\mathbf {x}) d \mu (\mathbf {x})
\]

对每个 \(\alpha\) 成立；而 \(\int h_{\alpha}(\mathbf{x}) d\mu (\mathbf{x}) = c_{\alpha} + \int \langle \mathbf{x}, \mathbf{z}_{\alpha}' \rangle d\mu (\mathbf{x})\)（因为 \(\mu\) 的总质量为 1）；但由重心的定义，\(\int \langle \mathbf{x}, \mathbf{z}_{\alpha}' \rangle d\mu (\mathbf{x}) = \langle \mathbf{b}_{\mu}, \mathbf{z}_{\alpha}' \rangle\)。因此 \(\sup \int h_{\alpha}(\mathbf{x}) d\mu (\mathbf{x}) = \sup h_{\alpha}(\mathbf{b}_{\mu}) = f(\mathbf{b}_{\mu})\)，结论得证。

当 \(\mu\) 是 K 上总质量为 1 的一个离散正测度时，命题 2 重新给出定义 K 上凸函数的那个不等式。

推论. --- 对每个凸数值函数 \(g\)（\(K\) 上的，有界且下半连续），有 \(g(\mathbf{b}_{\mu}) \leqslant \int g d\mu\) 。

只需注意到 \(\inf_{\mathbf{x} \in \mathrm{K}} g(\mathbf{x}) = a\) 有限，并对 \(g - a\) 应用命题 2。
% semantic-fragments: legacy-input-seam
\relax{}
